\section*{The Appointed Texts}

\begin{studybox}{Text boundary}
The following Scripture is the public-domain Douay--Rheims (Challoner) study
translation, not the English proclaimed at this celebration. The approved
U.S. Lectionary controls the proclamation. Protected Missal English is not
reproduced; Latin incipits and descriptions identify its proper units.
\end{studybox}

\subsection*{Exodus 16:2--4, 12--15}
All the congregation murmurs against Moses and Aaron, remembers Egypt's
fleshpots, and fears death by famine. The Lord promises to rain bread from
heaven, sufficient for each day, and says that the gift will test whether the
people walk in his law. Quails cover the camp at evening; dew surrounds it in
the morning; when the small frost-like thing appears, Israel asks what it is.
Moses answers: ``This is the bread which the Lord hath given you to eat.''

\subsection*{Psalm 78:3--4, 23--25, 54}
The assembly recalls what the ancestors told: God opened heaven's doors,
rained manna, gave the bread of heaven and abundant provision, and led the
people toward the sanctuary land. The appointed response is drawn from
v.~24b.

\subsection*{Ephesians 4:17, 20--24}
The Apostle forbids walking in the old vanity of mind. Those who have learned
Christ are to put off the old person, be renewed in the spirit of the mind,
and put on the new person created according to God in justice and holiness of
truth.

\subsection*{Matthew 4:4b and John 6:24--35}
The acclamation answers that one lives not by bread alone but by every word
from God. In the Gospel, the crowd crosses to Capharnaum seeking Jesus. He
exposes their attachment to the loaves, commands work for food enduring to
eternal life, and names belief in the one sent as God's work. When they invoke
the wilderness manna, Jesus identifies the Father as giver of the true bread
from heaven. He concludes: ``I am the bread of life. He that cometh to me
shall not hunger: and he that believeth in me shall never thirst.''

\subsection*{Missal units}
The Entrance begins \latin{Deus in loco sancto suo}; the Collect,
\latin{Protector in te sperantium, Deus}; the offerings prayer,
\latin{Suscipe, quaesumus, Domine, munera}; and the final prayer,
\latin{Sumpsimus, Domine, divinum sacramentum}. Communion offers Ps.~102:2
(modern 103) or John~6:35. Consult the approved Missal for the protected
English and the complete ritual text.
